% ============================================================
% Chương 7. Đánh giá  (chuẩn 2.2, 3.1, 3.2)
% Nguồn: docs/report_tv1.md mục 5, 6, 7; docs/optimization.md; docs/benchmark.md;
%        docs/data_models.md; docs/lo_traceability.md; tests/test_dashboard_questions.py
% ============================================================
\section{Đánh giá tổng thể}
\label{sec:danh-gia}

Hệ thống được đánh giá trên năm khía cạnh: độ đúng của kho qua bộ kiểm thử, sự nhất quán
giữa các hệ lưu trữ, khả năng trả lời các yêu cầu nghiệp vụ, hiệu quả của các kỹ thuật tối ưu
truy vấn, và so sánh hiệu năng giữa các cách tổ chức dữ liệu. Cuối chương đối chiếu kết quả
với yêu cầu của môn học và nêu các hạn chế.

% ------------------------------------------------------------
\subsection{Kiểm thử}
\label{sec:kiem-thu}

Bộ kiểm thử gồm 289 bài, viết bằng pytest. Kết quả lần chạy gần nhất: 286 bài đạt, 2 bài được
đánh dấu là lỗi đã biết (expected failure, trình bày ở mục~\ref{sec:doi-soat}), và 1 bài được bỏ
qua. Bài bị bỏ qua là bộ kiểm thử DAG, vì cần môi trường Airflow nên được chạy riêng trong
container Airflow, với 10 bài đều đạt. Các bài kiểm thử được chia theo nhóm như
Bảng~\ref{tab:nhom-kiem-thu}; số đúng của mỗi bài được chốt trước trong hợp đồng dữ liệu và
trong đặc tả yêu cầu nghiệp vụ, không lấy từ chính kết quả của pipeline.

\begin{table}[htbp]
\centering
\footnotesize
\begin{tabularx}{\textwidth}{|p{3.2cm}|X|r|}
\hline
\textbf{Nhóm} & \textbf{Nội dung kiểm tra} & \textbf{Số bài} \\ \hline
Số dòng & Bảng staging và 9 bảng kho đúng số dòng trong hợp đồng; PostgreSQL bằng Parquet & 25 \\ \hline
Cờ chất lượng & Số dòng \texttt{TRUE} của từng cờ đúng hợp đồng, trên Parquet và PostgreSQL & 75 \\ \hline
Khóa & Khóa chính không trùng; không có khóa ngoại mồ côi & 19 \\ \hline
Giá trị thiếu & Tỷ lệ \texttt{NULL} của từng cột đúng dự kiến; \texttt{NULL} được giữ nguyên khi nạp & 10 \\ \hline
Phụ thuộc hàm & Các phụ thuộc hàm làm căn cứ chuẩn hóa (mục~\ref{sec:chuan-hoa}) & 11 \\ \hline
Đối soát tổng & Tổng click, tổng điểm giữa nguồn và kho & 4 \\ \hline
Quy tắc nghiệp vụ & Cách tính tuần; bảng tiến độ đủ các tuần; cột lũy kế không giảm & 6 \\ \hline
Đáp án yêu cầu nghiệp vụ & Tỷ lệ trễ BR1; danh sách cảnh báo BR4; phân vị BR7 & 6 \\ \hline
Chống rò rỉ nhãn & Không có cột nhãn; đảo kết quả cuối kỳ không đổi danh sách BR4 & 2 \\ \hline
Chạy lại và stream & Nạp lại không nhân đôi dữ liệu; kết quả stream khớp batch & 5 \\ \hline
Kho mở rộng 5$\times$ & Số dòng gấp đúng 5 lần; tỷ lệ và phân bố không đổi & 19 \\ \hline
Báo lỗi Airflow & Ghi log, gửi webhook đúng định dạng; webhook lỗi không làm hỏng lần chạy & 5 \\ \hline
Câu hỏi mẫu (SQL) & 20 đáp án và bài đảo nhãn của câu 14 & 21 \\ \hline
Câu hỏi mẫu (dashboard) & Tìm được đáp án của 20 câu trên giao diện dashboard & 20 \\ \hline
Dashboard & Các trang khởi động và hiển thị không lỗi & 4 \\ \hline
Truy vấn, data mart, benchmark & SQL nghiệp vụ, hai data mart, các trang và bộ chạy benchmark & 56 \\ \hline
DAG Airflow & Đủ task, đúng thứ tự, cổng chất lượng đứng trước hai bước nạp & 1 (chạy riêng: 10) \\ \hline
\textbf{Tổng} & & \textbf{289} \\ \hline
\end{tabularx}
\caption{Phân nhóm bộ kiểm thử}
\label{tab:nhom-kiem-thu}
\end{table}
% Lưu ý: dòng "Truy vấn, data mart, benchmark" gộp 5 bài SQL của TV2 và 51 bài của TV3.
% Nếu số bài thay đổi, cập nhật lại theo `make test`.

% ------------------------------------------------------------
\subsection{Đối soát và độ đúng của câu hỏi mẫu}
\label{sec:doi-soat}

Sau mỗi lần nạp, số dòng của cả 9 bảng trên BigQuery, PostgreSQL và Parquet khớp nhau. Ở
mức nghiệp vụ, SQL của 20 câu hỏi mẫu chạy trên cả ba hệ đều cho đáp án trùng với đáp án
chuẩn, vốn được tính trực tiếp từ 7 file CSV mà không đi qua kho. Kết quả này cho thấy quá
trình ETL không làm sai lệch dữ liệu và ba hệ lưu trữ nhất quán với nhau.

Ngoài việc kiểm tra SQL, nhóm còn kiểm tra người dùng có tìm được đáp án trên giao diện hay
không: mỗi câu hỏi được mô phỏng bằng cách mở đúng trang, chọn đúng đợt học và tuần như đề
hỏi, rồi tìm đáp án trong những gì trang hiển thị. Dashboard trả lời được 18 trên 20 câu. Hai
câu chưa trả lời được là câu 7 (tỷ lệ nộp trễ chung của mọi bài TMA) và câu 8 (điểm TMA trung
bình theo môn), vì trang Nộp bài và điểm hiện chỉ cho xem từng đợt học riêng lẻ. Ngoài ra, câu
15 và câu 16 tuy đạt nhưng đáp án đang được ghi cố định trong trang Cảnh báo sớm và rút lui,
chưa được tính từ dữ liệu.

% ------------------------------------------------------------
\subsection{Kết quả các yêu cầu nghiệp vụ}
\label{sec:ket-qua-br}

Bảng~\ref{tab:ket-qua-br} tóm tắt kết quả tiêu biểu của từng yêu cầu nghiệp vụ trên toàn bộ
dữ liệu hoặc trên đợt học được chọn làm ví dụ trong đặc tả.

\begin{table}[htbp]
\centering
\small
\begin{tabularx}{\textwidth}{|l|X|}
\hline
\textbf{BR} & \textbf{Kết quả tiêu biểu} \\ \hline
BR1 & 16,8\% bài TMA nộp trễ (16.320 bài) \\ \hline
BR2 & \texttt{oucontent} chiếm 28,3\% tổng click; tuần 2 có nhiều click nhất (2.307.080) \\ \hline
BR3 & Cả 19 đợt đều có $\rho$ dương, từ 0,16 (BBB 2014J) đến 0,41 (EEE 2014J): liên hệ dương nhưng yếu đến vừa \\ \hline
BR4 & Tuần 4: 3.958 sinh viên vào danh sách; 80,6\% trong số đó cuối cùng rớt hoặc rút, so với tỷ lệ nền 43,5\%; độ bao phủ 26,9\% \\ \hline
BR5 & 30,4\% số lượt rút diễn ra từ ngày khai giảng trở về trước; ngày rút trung vị là ngày 27 \\ \hline
BR6 & Kỳ 2014J: nhóm IMD 90--100\% đạt 63,3\%, nhóm 0--10\% đạt 41,1\%, chênh 22,2 điểm phần trăm (kỳ 2013B chênh 25,3 điểm) \\ \hline
BR7 & AAA 2014J, tuần 10: trung vị tổng click của 336 sinh viên còn học là 593 \\ \hline
BR8 & BBB 2013J, tuần 5: nhóm đạt (1.072 sinh viên) có trung vị 5 click vào diễn đàn, khoảng 25\%--75\% là 0--22 \\ \hline
\end{tabularx}
\caption{Kết quả tiêu biểu của 8 yêu cầu nghiệp vụ}
\label{tab:ket-qua-br}
\end{table}

Kết quả đáng chú ý nhất là quy tắc cảnh báo sớm của BR4. Chỉ với dữ liệu đến hết tuần 4, quy
tắc 4 dấu hiệu chọn ra 3.958 sinh viên, trong đó 80,6\% thực sự rớt hoặc rút khi đợt học kết
thúc, gần gấp đôi tỷ lệ nền 43,5\%. Độ bao phủ 26,9\% cho thấy quy tắc ưu tiên độ chính xác:
danh sách ngắn nhưng phần lớn là đúng, phù hợp khi nguồn lực hỗ trợ sinh viên có hạn. Kết quả
của BR3 cần được đọc thận trọng: hệ số tương quan dương ở mọi đợt chỉ cho thấy liên hệ, không
cho phép kết luận rằng hoạt động nhiều hơn dẫn đến điểm cao hơn.

% TODO: mỗi yêu cầu (hoặc mỗi nhóm người dùng) một ảnh chụp dashboard minh họa.

% ------------------------------------------------------------
\subsection{Tối ưu truy vấn}
\label{sec:toi-uu}

Bốn kỹ thuật tối ưu được thử nghiệm trên các truy vấn thực tế của dashboard
(Bảng~\ref{tab:toi-uu}). Mỗi truy vấn chạy 5 lần, lấy trung vị; bộ nhớ đệm kết quả của
BigQuery được tắt. Với mỗi kỹ thuật, kết quả trước và sau tối ưu được so sánh và phải giống
hệt nhau. Trên BigQuery, chỉ số chính là lượng dữ liệu quét, vì chi phí và hạn mức được tính
theo chỉ số này; ở quy mô dưới 1 GB, thời gian chạy chủ yếu là chi phí khởi động job nên
chênh lệch không đáng kể.

\begin{table}[htbp]
\centering
\small
\begin{tabularx}{\textwidth}{|p{3.6cm}|X|l|r|r|r|}
\hline
\textbf{Kỹ thuật} & \textbf{Truy vấn} & \textbf{Chỉ số} & \textbf{Trước} & \textbf{Sau} & \textbf{Giảm} \\ \hline
Phân vùng bảng tiến độ theo tuần (BigQuery) & Danh sách cảnh báo BR4 tuần 4 & dữ liệu quét & 79,2 MB & 4,7 MB & 16,8$\times$ \\ \hline
Gom cụm bảng click theo đợt (BigQuery) & Click theo tuần và loại học liệu của một đợt & dữ liệu quét & 304,6 MB & 12,0 MB & 25,5$\times$ \\ \hline
Bảng tổng hợp sẵn BR2 (BigQuery) & Biểu đồ BR2 của một đợt & dữ liệu quét & 381,6 MB & 1,2 MB & 330$\times$ \\ \hline
Chỉ mục B-tree (PostgreSQL) & Như kỹ thuật gom cụm & thời gian & 400,1 ms & 73,6 ms & 5,4$\times$ \\ \hline
\end{tabularx}
\caption{Kết quả các kỹ thuật tối ưu truy vấn}
\label{tab:toi-uu}
\end{table}

Các phép đo cho thấy mỗi kỹ thuật chỉ có tác dụng trong điều kiện nhất định. Phân vùng chỉ
giảm dữ liệu quét khi truy vấn lọc cột phân vùng bằng một hằng số: truy vấn BR4 nối bảng với
chính nó qua điều kiện tuần trước bằng tuần hiện tại trừ một, và BigQuery không dùng được
điều kiện này để bỏ qua phân vùng; chỉ khi thêm điều kiện hằng số tương đương thì dữ liệu quét
mới giảm từ 79,2 MB xuống 4,7 MB. Gom cụm phù hợp với các cột chuỗi hay được lọc như mã môn
và mã kỳ, khi đó truy vấn một đợt chỉ đọc khoảng 1/22 bảng. Bảng tổng hợp sẵn đổi chi phí của
mỗi lần mở trang lấy chi phí dựng bảng một lần, nhưng phải được dựng lại sau mỗi lần nạp kho.
Chỉ mục trên PostgreSQL phải bắt đầu bằng cột được lọc: khóa chính có sẵn bắt đầu bằng mã
sinh viên nên không giúp được truy vấn theo đợt; chỉ mục mới kèm các cột cần đọc cho phép
truy vấn chỉ đọc chỉ mục mà không mở bảng, đổi lại chỉ mục chiếm 499 MB và làm chậm việc nạp.

% ------------------------------------------------------------
\subsection{So sánh các cách tổ chức dữ liệu}
\label{sec:benchmark}

Ba cách tổ chức dữ liệu được so sánh: dimensional model (star) của kho chính, lược đồ 3NF
(mục~\ref{sec:chuan-hoa}) và bảng phi chuẩn hóa toàn phần (wide table), trong đó mỗi bảng fact
được nối sẵn mọi thuộc tính của dimension. Bốn truy vấn đại diện được chọn từ các câu hỏi mẫu:
BR1 (câu 7), BR2 (câu 9), BR4 (câu 14) và BR6 (câu 18). Mỗi tổ hợp truy vấn, hệ quản trị
(PostgreSQL, BigQuery) và quy mô (1$\times$, 5$\times$) được chạy 5 lần, lấy trung vị. Có 44
trên 48 tổ hợp được đo; bốn tổ hợp wide table 5$\times$ trên BigQuery bị loại vì vượt giới hạn
10 GB của gói Sandbox. Số liệu dưới đây là của lần đo ngày 26/10/2026.

\begin{table}[htbp]
\centering
\small
\begin{tabular}{|l|l|r|r|r|}
\hline
\textbf{Truy vấn} & \textbf{Hệ, quy mô} & \textbf{Star (s)} & \textbf{3NF} & \textbf{Wide table} \\ \hline
BR1 (câu 7) & PostgreSQL 5$\times$ & 0,047 & 1,12$\times$ & 1,18$\times$ \\ \hline
BR2 (câu 9) & PostgreSQL 5$\times$ & 3,418 & 1,16$\times$ & 1,56$\times$ \\ \hline
BR6 (câu 18) & PostgreSQL 5$\times$ & 0,036 & 0,91$\times$ & 0,67$\times$ \\ \hline
BR2 (câu 9) & BigQuery 5$\times$ & 1,612 & 0,40$\times$ & -- \\ \hline
BR4 (câu 14) & BigQuery 5$\times$ & 3,098 & 0,22$\times$ & -- \\ \hline
\end{tabular}
\caption{Thời gian trung vị của star và tỷ lệ so với star (lớn hơn 1 là chậm hơn), một số tổ hợp tiêu biểu}
\label{tab:benchmark}
\end{table}

Trên PostgreSQL, hệ lưu trữ theo dòng, ba cách tổ chức chênh nhau không quá 16\% ở quy mô
1$\times$, vì toàn bộ dữ liệu nằm vừa bộ nhớ và các phép nối với dimension nhỏ đều rẻ. Ở quy mô
5$\times$, wide table chậm nhất ở truy vấn quét bảng click (BR2 chậm hơn star 1,56 lần), vì
bảng rộng gấp đôi star và đọc bảng là đọc cả dòng, kể cả các cột không dùng. Lược đồ 3NF chậm
hơn star 8--16\% ở BR1 và BR2, vì phải nối lại dimension để khôi phục mã đợt. Ngược lại, wide
table nhanh hơn ở truy vấn nhỏ có phép nối như BR6.

Trên BigQuery, hệ lưu trữ theo cột, star và 3NF quét đúng cùng một lượng dữ liệu ở cả bốn truy
vấn, vì truy vấn chỉ đọc các cột cần thiết và các cột dư thừa của star không bị đọc. Tuy nhiên,
ở quy mô 5$\times$, star chậm hơn 3NF 2,5 lần (BR2) và 4,5 lần (BR4) về thời gian dù quét cùng
dữ liệu. Khác biệt duy nhất giữa hai trường hợp là cách tạo bảng: kho star được nạp bằng load
job từ file Parquet, còn bảng 3NF được tạo bằng câu lệnh SQL chạy phân tán. Đây là giải thích
hợp lý nhất nhưng chưa được kiểm chứng trực tiếp, vì BigQuery không cho xem cách bảng được chia
khối lưu trữ.

Không cách tổ chức nào nhanh nhất ở mọi tổ hợp. Nhóm giữ star làm kho chính vì star dễ viết
truy vấn nhất, nhanh hơn 3NF trên PostgreSQL và quét cùng lượng dữ liệu trên BigQuery. Việc
tăng tốc cho từng trang dashboard được thực hiện bằng phân vùng, gom cụm, chỉ mục và bảng tổng
hợp sẵn (mục~\ref{sec:toi-uu}) thay vì thay đổi cả cách tổ chức.

% ------------------------------------------------------------
\subsection{Đối chiếu với yêu cầu môn học}
\label{sec:doi-chieu}

Bảng~\ref{tab:tu-danh-gia} tóm tắt nội dung của project ứng với từng chuẩn đầu ra của môn học.

\begin{table}[htbp]
\centering
\small
\begin{tabularx}{\textwidth}{|p{3.0cm}|X|X|l|}
\hline
\textbf{Chuẩn đầu ra} & \textbf{Nội dung trong project} & \textbf{Con số tiêu biểu} & \textbf{Mục} \\ \hline
1.1 Thiết kế CSDL ở 3 mức & Thiết kế ý niệm, luận lý, vật lý; kiểm chứng phụ thuộc hàm trên dữ liệu & 8 thực thể $\rightarrow$ 5 dimension và 4 fact; 11 phụ thuộc hàm được kiểm thử & \ref{sec:thiet-ke} \\ \hline
1.2 Mô hình hóa từ quan hệ đến phi quan hệ & 6 mô hình dữ liệu, có số đo so sánh & Dạng lồng quét ít hơn 33\%; JSONB nhỏ hơn 7,4$\times$ & \ref{sec:phi-quan-he} \\ \hline
2.1 Kiến trúc quản lý dữ liệu từ nhiều nguồn & 4 tầng dữ liệu, 3 hệ lưu trữ, điều phối bằng Airflow & Dựng lại toàn bộ kho trong 303 giây & \ref{sec:kien-truc}, \ref{sec:hien-thuc} \\ \hline
2.2 Thiết kế kho dữ liệu & Lược đồ constellation, bảng ảnh chụp theo tuần, so sánh 3 cách tổ chức & 1.342.949 dòng ảnh chụp; 44 tổ hợp đo hiệu năng & \ref{sec:luan-ly}, \ref{sec:benchmark} \\ \hline
3.1 Quản lý dữ liệu phức tạp & Gộp dòng trùng, cờ chất lượng, chống rò rỉ nhãn, đối soát ba hệ & 3.810.465 dòng trùng được gộp; 289 bài kiểm thử & \ref{sec:etl}, \ref{sec:kiem-thu} \\ \hline
3.2 Xử lý dữ liệu theo kiểu mẫu hiện đại & Xử lý phân tán (Spark), kho đám mây (BigQuery), điều phối (Airflow), tối ưu có số đo & Giảm 5,4$\times$ đến 330$\times$ & \ref{sec:hien-thuc}, \ref{sec:toi-uu} \\ \hline
% TODO (bổ sung sau): thêm bằng chứng stream vào 1.2, 2.1, 3.2
\end{tabularx}
\caption{Nội dung project ứng với chuẩn đầu ra của môn học}
\label{tab:tu-danh-gia}
\end{table}

% ------------------------------------------------------------
\subsection{Hạn chế}
\label{sec:han-che}

Hệ thống còn một số hạn chế. Về dữ liệu, 7 nguồn đều thuộc cùng một bộ dữ liệu tĩnh, nên bài
toán tích hợp nhiều nguồn mới thể hiện ở sự khác biệt về khóa, chất lượng và quy ước ngày giữa
các file, chưa có các nguồn thuộc hệ thống khác nhau. Bản mở rộng quy mô 5$\times$ được tạo
bằng cách nhân bản, nên chỉ phù hợp để đo tải, không bổ sung sự đa dạng thật của dữ liệu.

Về hạ tầng, gói Sandbox của BigQuery tự xóa bảng sau 60 ngày và không cho dùng các lệnh sửa
dữ liệu, nên kho được nạp lại toàn bộ mỗi lần thay vì nạp tăng dần, và hệ thống chưa có cơ chế
bắt thay đổi dữ liệu (CDC). Cũng vì giới hạn dung lượng 10 GB, wide table ở quy mô 5$\times$
không dựng được trên BigQuery.

Về hiện thực, các kỹ thuật phân vùng, gom cụm và chỉ mục mới được thử nghiệm trên bản sao của
bảng, chưa được đưa vào luồng nạp chính. Hai data mart được làm mới thủ công, nằm ngoài DAG, nên
có thể hiển thị số cũ nếu kho được nạp lại mà quên làm mới. Ba trang dashboard của BR1, BR4,
BR5 và BR7 vẫn đọc trực tiếp Parquet thay vì đi qua lớp truy vấn chung, và dashboard chưa trả
lời được 2 trong 20 câu hỏi mẫu (mục~\ref{sec:doi-soat}). Cuối cùng, kết quả so sánh star với
3NF trên BigQuery ở quy mô 5$\times$ còn một điểm chưa giải thích được bằng thực nghiệm
(mục~\ref{sec:benchmark}).
